\section{From the Proposal to the Current Project}
\label{sec:proposal}

The proposal for this project asked whether a server that streams a changing 3D scene to a web browser should schedule its updates by the capacity the browser measures and reports. % the submitted title is not recorded in the repo; the form text is in the 2026-11-02 session
It placed the system between a server and a web browser.
A server would stream 3DGS scene updates over WebSocket to a renderer written in WebGPU, as additions, modifications and deletions of Gaussians, each addressed by a Gaussian ID.
The browser would keep a fixed pool of GPU slots and a table that maps Gaussian IDs to slots, in the way an operating system keeps page frames and a page table, so that an addition fills a free slot and a deletion releases one.

The plan was to measure before designing.
It would first measure the renderer's frame time against scene size, the highest update rate the browser could absorb without dropping frames, the latency of each stage of the pipeline, and the pipeline's behavior under emulated bandwidth limits.
It would then compare sender policies: first-in first-out, coalescing to the latest value of each Gaussian, flow control by acknowledgements, and admission driven by the capacity the client reports.
The proposal likened the last policy to the balloon driver through which a hypervisor reclaims memory from a guest with the guest's cooperation~\cite{waldspurger2002}.
Two clients sharing one network bottleneck would test whether one client's load disturbs the other, the noisy-neighbor problem of shared hosts.
Its related work covered streaming Gaussian scenes to clients and rendering them in the browser~\cite{siekkinen2026,websplatter2026,moqsplat2026,lts2025,lmg2026,gaussanything2026,partialgs2026}, with 3DGS as the representation~\cite{kerbl2023gaussian} and StreamSplat as a candidate source of updates~\cite{streamsplat2026}.
The reconstruction model was deferred to the last milestone, so that the streaming pipeline could be built and understood first.

The virtualization in that plan was an analogy.
Slots and page tables described how a renderer manages its own buffers, and no part of the system shared a physical device among tenants through a virtualization mechanism.
Real GPU sharing appeared only as an optional extension that would run several model tenants on the server GPU under time-slicing and MPS; it needed the model, so it would come after the model milestone or use synthetic GPU load as its tenants, and it was the first milestone to be cut if time ran short.
On the day the proposal was submitted, we therefore moved the virtualization core of the project to the server.
In the current project many users stream live video, one GPU runs a reconstruction per user, the model is part of the system from the start, and the results are to stream back to each user's browser.
The question changed accordingly, from how a sender should pace updates to one client to how many users one GPU can serve and which lever raises that number.
The browser part of the original plan is now future work (Section~\ref{sec:remaining}).

Before settling on a testbed we probed which sharing mechanisms the available infrastructure exposes.
These probes are described here for context and are not measurements in this report.
On the H200 nodes of the NCHC nano4 cluster MIG was disabled, so hardware partitions were not available.
A user-started MPS daemon worked inside a batch job, and its per-client memory limit applied to clients started after the limit was set, while a client that was already running could not grow after the limit was raised.
Inside a Slurm job there was no cgroup delegation, so a rootless Kubernetes installation would have no way to enforce resource limits on its containers.
On our local workstation, a kind cluster with the NVIDIA device plugin scheduled a pod onto the RTX 6000 Ada, which in the plugin's default configuration gives the pod the whole GPU~\cite{k8s_device_plugin}.
All measurements in this report run on that local RTX 6000 Ada, with each user served either by its own process or container under time-slicing or MPS, or by one server process shared by all users (Section~\ref{sec:setup}).
